\tocsection{\textbf{9/14/26 Lecture 8:} Product Spaces}

\textbf{Question:} Given a collection of subsets of $X$, $\beta$, when is it the basis for some topology?

\begin{lemma}[\textbf{Lemma 8.1}]
    A collection $\beta$ of subsets is the basis for some topology on $X$ if, and only if $X = \bigcup_{B \in \beta} B$, and finite intersections of elements in $\beta$ is a union of elements in $\beta$. 
\end{lemma}
\begin{proof}{Proof:}
    $(\Rightarrow)$ This follows from the definition of a basis.\\

    $(\Leftarrow)$ Suppose that $\beta$ is a collection of subsets that have the above property, and say that $U$ is open iff it is the union of elements in $\beta$. We will show that this is a topology.\\

    First, we see that $\emptyset$ and $X$ are open. Now, if $\{A_{\lambda}\}$ are open, then we see that $\bigcup_{\lambda} A_{\lambda}$ is open, because each of the $A_{\lambda}$ are unions of elements in $\beta$, so $\bigcup_{\lambda} A_{\lambda}$ is too.\\

    Finally, if $A_1, A_2$ are open, with $A_1 = \bigcup_{B_{\lambda, 1} \in \beta} B_{\lambda, 1}$ and $A_2 = \bigcup_{B_{\lambda, 2} \in \beta} B_{\lambda, 2}$, then we see that
    \[ A_1 \cap A_2 = \bigcup_{\lambda_1, \lambda_2} (B_{\lambda, 1} \cap B_{\lambda, 2}) \]
    and as each of the $B_{\lambda, 1}, B_{\lambda, 2}$ are unionts of elements in $\beta$, we see that so is $A_1 \cap A_2$. Thus, by induction, all finite intersections are as well.

    $\qed$
\end{proof}

\subsection*{Product Spaces}

Let $X, Y$ be topological spaces. We will consider $X \times Y = \{(x, y) \col x \in X, y \in Y\}$ and see what natural topologies we can put on it.\\

\textbf{First Attempt:} What if we just let the open sets be of the form $\cal{U} \times \cal{V}$ for open sets $\cal{U}, \cal{V}$ of $X$ and $Y$ respectively?

\textbf{Problem:} Quite easy to show that this is not closed under unions.\\

\textbf{Second Attempt:} How about, in the fashion of Lemma 8.1, we define the basis of our topology to be $\beta = \{\cal{U} \times \cal{V} \col \cal{U} \text{ open in }X, \cal{V} \text{ open in }Y\}$. It turns out this works, because
\begin{enumerate}
    \item $X \times Y$ is in $\beta$
    \item $(U_1 \times V_1) \cap (U_2 \times V_2) = (U_1 \cap U_2) \times (V_1 \cap V_2)$, so intersections of basis elements are basis elements!
\end{enumerate}
Thus, we now have our naturally arising topology.

\begin{definition}[\textbf{Definition 8.2:} Product Spaces]
    Let $X, Y$ be topological spaces. The \emph{product space} is the set $X \times Y$ equipped with the topology generated by the basis $\beta = \{\cal{U} \times \cal{V} \col \cal{U} \text{ open in }X, \cal{V} \text{ open in }Y\}$, called the \emph{product topology}.
\end{definition}

\begin{definition}[\textbf{Definition 8.3:} Projections]
    The maps $p_1 : X \times Y \to X$ and $p_2 : X \times Y \to Y$, called \emph{projections} are given by
    \[ p_1(x, y) = x \quand p_2(x, y) = y \]
    These are continuous and are also open maps (meaning the image of an open set is open).
\end{definition}
The proof of the claims made above are quite simple. First, just consider $p_1$. Then, $p_1^{-1}(U) = U \times Y$, so if $U$ is open, the preimage is open as well. Similarly, if $\bigcup_{\alpha} U_{\alpha} \times V_{\alpha}$ is open in $X \times Y$, then we see that
\[ p_1\left(\bigcup_{\alpha} U_{\alpha} \times V_{\alpha}\right) = \bigcup_{\alpha} U_{\alpha} \]
which is open.

\begin{lemma}[\textbf{Proposition 8.4}]
    A function $f : Z \to X \times Y$ if, and only if $p_1 \circ f : Z \to X$ and $p_2 \circ f : Z \to Y$ are continuous.
\end{lemma}
\begin{proof}{Proof:}
    First, as continuity is preserved under composition, we see that if $f$ is continuous, so are $p_1 \circ f$ and $p_2 \circ f$.\\

    Now, suppose that $p_1 \circ f$ and $p_2 \circ f$ are continuous. Then, as
    \[ f^{-1}\left(\bigcup U_{\alpha} \times V_{\alpha}\right) = \bigcup f^{-1}(U_{\alpha} \times V_{\alpha}) \]
    it suffices to show that the preimage of basis elements are open. Well, we see that
    \[ p_1^{-1}(U_{\alpha}) = U_{\alpha} \times Y \quand p_{2}^{-1}(V_{\alpha}) = X \times V_{\alpha} \]
    so
    \[ U_{\alpha} \times V_{\alpha} = p_1^{-1}(U_{\alpha}) \cap p_2^{-1}(V_{\alpha}) \]
    Thus, 
    \begin{align*}
        f^{-1}(U_{\alpha} \times V_{\alpha}) &= f^{-1}(p_1^{-1}(U_{\alpha}) \cap p_{2}^{-1}(V_{\alpha}))\\
        &= (p_1 \circ f)^{-1}(U_{\alpha}) \cap (p_2 \circ f)^{-1}(V_{\alpha})
    \end{align*}
    which is open. 

    $\qed$
\end{proof}

\begin{lemma}[\textbf{Proposition 8.5}]
    $X \times Y$ is Hausdorff if, and only if, $X$ and $Y$ are Hausdorff. 
\end{lemma}
\begin{proof}{Proof:}
    $(\Leftarrow)$ Suppose that $X$ and $Y$ are Hausdorff, and let $(x_1, y_1) \ne (x_2, y_2)$. Then, WLOG, assume that $x_1 \ne x_2$. Then, let $U_1$ and $U_2$ be disjoint open neighborhoods of $X$ containing $x_1$ and $x_2$ respectively. Then, we see that $U_1 \times Y$ and $U_2 \times Y$ are disjoint open neighborhoods containing $(x_1, y_1)$ and $(x_2, y_2)$ respectively. Thus, $X \times Y$ is Hausdorff.\\

    $(\Rightarrow)$ Suppose that $X \times Y$ is Hausdorff. We will show that $X$ is Hausdorff, and the proof for $Y$ follows similarly. Let $x_1 \ne x_2$. If $Y$ is empty, then $X \times Y = X$, so it follows. Otherwise, pick a $y \in Y$, and consider $(x_1, y)$ and $(x_2, y)$. These are nonequal, so let $(U_1 \times V_1)$ and $(U_2 \times V_2)$ be disjoint open basis elements containing $(x_1, y)$ and $(x_2, y)$ respectively (recall that we can always pass open neighborhoods to basis elements). Then, notice that
    \[ x_1 \in U_1, x_2 \in U_2 \]
    and $U_1$ and $U_2$ are open. Also, if some $p$ lies in both, then $(p, y)$ lies in both $U_1 \times V_1$ and $U_2 \times V_2$, which is a contradiction. Thus, they are disjoint open neighborhoods, so $X$ is Hausdorff. 

    $\qed$
\end{proof}

\begin{lemma}[\textbf{Lemma 8.6}]
    If $X$ is a topological space and $\beta$ is a basis, then $X$ is compact if, and only if, every covering of $X$ by basis elements has a finite subcover. 
\end{lemma}
\begin{proof}{Proof:}
    $(\Rightarrow)$ Obvious.\\

    $(\Leftarrow)$ Suppose that every covering of $X$ by basis elements has a finite subcover. Now, let $\{U_{\lambda}\}$ be an open cover of $X$. Then, we can write $U_{\lambda} = \bigcup_{\alpha} B_{\alpha, \lambda}$ for some $B_{\alpha, \lambda} \in \beta$. Thus, as
    \[ X = \bigcup_{\lambda, \alpha} B_{\alpha, \lambda} \]
    we can pass to a finite subcover
    \[ X = \bigcup_{i = 1}^{n} B_{\alpha_i, \lambda_i} \]
    Then, as $B_{\alpha_i, \lambda_i} \subseteq U_{\lambda_i}$, we see that
    \[ X = \bigcup_{i = 1}^{n} U_{\lambda_i} \]
    which is a finite subcover.

    $\qed$
\end{proof}